\section{Conclusões}\label{sec:conclusoes}

Este trabalho apresentou uma abordagem de projeto inteligente para o dimensionamento de pontes de madeira roliça, integrando modelagem paramétrica, verificação normativa segundo a ABNT NBR 7190 e otimização robusta multiobjetivo por meio do algoritmo NSGA-II. Os resultados demonstram que os objetivos propostos foram atingidos, tanto do ponto de vista computacional quanto estrutural.

A plataforma computacional desenvolvida mostrou-se eficiente e funcional, permitindo a obtenção de soluções de projeto em curto intervalo de tempo por meio de uma interface intuitiva. Essa característica representa um avanço em relação aos métodos tradicionais de dimensionamento, ao possibilitar a geração rápida de múltiplas alternativas e favorecer decisões técnicas mais fundamentadas.

O algoritmo NSGA-II demonstrou elevada eficiência na obtenção da fronteira de Pareto, identificando soluções estruturalmente viáveis e materialmente eficientes. A comparação com o método de Monte Carlo evidenciou a superioridade do processo de otimização, uma vez que as soluções obtidas por amostragem aleatória apresentaram comportamento mais conservador, com maior consumo de material e baixa utilização da capacidade resistente da estrutura.

A incorporação da robustez, avaliando cada solução candidata sob um desvio $\rho$ aplicado aos valores nominais das variáveis de projeto, mostrou-se essencial para garantir que as soluções da fronteira eficiente permaneçam seguras diante da variabilidade dimensional inerente à madeira roliça. Ao penalizar soluções posicionadas exatamente sobre as restrições ativas, a estratégia deslocou a fronteira para o interior do domínio viável, assegurando margem de segurança adicional em troca de um acréscimo controlado no consumo de material. O estudo de sensibilidade do parâmetro de robustez ($\rho = 0\%,\ 2{,}5\%,\ 5\%,\ 10\%$) quantificou esse acréscimo -- o custo da robustez --, \tofill{[A PREENCHER: sintetizar aqui a faixa de acréscimo percentual de volume obtida na Seção~\ref{sec:custo_robustez}]}, e a validação por simulação de Monte Carlo independente confirmou empiricamente a redução da probabilidade de violação das restrições na solução robusta frente à determinística de mesmo volume (Seção~\ref{sec:validacao_mc}).

A análise comparativa entre os cenários de carregamento confirmou a coerência física do modelo: a Ponte~02, com vão maior e trem tipo mais severo (TB-450), demandou \tofill{[A PREENCHER: acréscimo \% de volume, Seção~\ref{sec:comparacao}]} a mais de madeira do que a Ponte~01, resultado consistente com a elevação combinada do vão e das solicitações de projeto. A análise de sensibilidade global (índices de Sobol) e a caracterização da dispersão das variáveis de projeto ao longo da fronteira eficiente (Seções~\ref{sec:sobol} e~\ref{sec:dispersao}) indicaram que \tofill{[A PREENCHER: variável(is) de maior influência sobre as restrições e sobre a diversidade da fronteira]}, informação diretamente aplicável ao controle de qualidade na fabricação e na montagem das peças de madeira roliça.

Como perspectivas de continuidade, destaca-se a integração ao modelo dos demais componentes do sistema estrutural, como ligações, sistemas de fundação e pilares, cuja influência é direta no comportamento global da estrutura. Recomenda-se ainda a comparação entre a estratégia de agregação por média e por pior caso ($\bar{g}_j = \max_c g_j$), a conexão com análises de confiabilidade estrutural baseadas em FORM ou Monte Carlo, e a validação da metodologia frente a um estudo de caso com ponte real de estrada vicinal, ampliando o alcance da metodologia proposta.
